\subsection{Boundedification}

Let E be a complex Hilbert space. We denote by s the unitary map \((x,y)\mapsto(-y,x)\) on \(E\oplus E\) . Let u be a closed partial operator with dense domain on E. The partial operator \(u^{*}u\) is self-adjoint and positive (prop. 12 of IV, p. 241), hence \(-1\notin\mathrm{Sp}(u^{*}u)\) (prop. 17 of IV, p. 248). We denote by \(\mathrm{W}(u)=(1_{\mathrm{E}}+u^{*}u)^{-1}=-\mathrm{R}(u^{*}u,-1)\) ; this is a positive injective endomorphism of E.

Denote by \(p_{1}\) and \(p_{2}\) the two canonical projections of \(\Gamma_{u}\) into E. These are elements of \(\mathcal{L}(\Gamma_{u};\mathrm{E})\) , and we have the equality of correspondences \(p_{2}=u\circ p_{1}\) . Let \((j,|p_{1}^{*}|)\) be the polar decomposition (def. 4 of I, p. 140) of the endomorphism \(p_{1}^{*}\in\mathcal{L}(\mathrm{E};\Gamma_{u})\) , so that \(p_{1}^{*}=|p_{1}^{*}|\circ j\) . The map j is a partial isometry from E into \(\Gamma_{u}\) .

DEFINITION 1. --- The endomorphism \(p_2 \circ j\) of E is called the boundedification of \(u\) ; we denote by \(b(u)\) this endomorphism.

One has \(\|b(u)\|\leqslant1\) .

PROPOSITION 1. --- We have the formulas

\[
| p _ {1} ^ {*} | = \mathrm{W} (u) ^ {1 / 2},\tag{1}
\]

\[
b (u) = u \circ \mathrm{W} (u) ^ {1 / 2} = u \circ | p _ {1} ^ {*} |,\tag{2}
\]

\[
1 _ {\mathrm{E}} - b (u) ^ {*} b (u) = \mathrm{W} (u),\tag{3}
\]

\[
b (u) \mathrm{W} (u) = \mathrm{W} (u ^ {*}) b (u), \quad b (u) \mathrm{W} (u) ^ {1 / 2} = \mathrm{W} (u ^ {*}) ^ {1 / 2} b (u),\tag{4}
\]

Let \(x \in E\) . Set \(y = \mathrm{W}(u)(x) \in \mathrm{dom}(u^{*}u)\) . We have \(y \in \mathrm{dom}(u)\) and \(u(y) \in \mathrm{dom}(u^{*})\) . For every element \((y_{1}, u(y_{1}))\) of \(\Gamma_{u}\) , where \(y_{1} \in \mathrm{dom}(u)\) , one computes

\[
\begin{array}{c} \langle (y, u (y)) | (y _ {1}, u (y _ {1})) \rangle = \langle y | y _ {1} \rangle + \langle u (y) | u (y _ {1}) \rangle \\ = \langle y + u ^ {*} u (y) | y _ {1} \rangle = \langle x | p _ {1} (y _ {1}, u (y _ {1})) \rangle . \end{array}
\]

This means that \(p_{1}^{*}(x) = (y, u(y))\) , whence \(p_{1} \circ p_{1}^{*}(x) = y = \mathrm{W}(u)(x)\) . Consequently, one has \(p_{1} \circ p_{1}^{*} = \mathrm{W}(u)\) , hence \(|p_{1}^{*}| = \mathrm{W}(u)^{1/2}\) , that is, formula (1). In particular, since \(\operatorname{Im}(p_{1}) = \operatorname{Im}(|p_{1}^{*}|)\) by the cor. of prop. 11 of I, p. 140, it follows that \(\operatorname{dom}(u) = \operatorname{Im}(p_{1}) = \operatorname{Im}(\mathrm{W}(u)^{1/2})\) .

The partial operator \(u \circ \mathrm{W}(u)^{1/2}\) therefore has domain E. The relation \(p_1 \circ j = (p_1^*)^* \circ j = |p_1^*| = \mathrm{W}(u)^{1/2}\) (prop. 11, a) of I, p. 140) then implies that

\[
b (u) = p _ {2} \circ j = u \circ p _ {1} \circ j = u \circ \mathrm{W} (u) ^ {1 / 2},
\]

which is formula (2).

One has \(\operatorname{Ker}(j) = \operatorname{Ker}(|p_1^*|)\) (prop. 10, \(b\) ) of I, p. 139) and formula (1) therefore imply that the partial isometry \(j: \mathrm{E} \to \Gamma_u\) is injective, hence isometric, whence \(j^* \circ j = 1_{\mathrm{E}}\) . Writing \(j = (p_1 \circ j, b(u))\) , we find

\[
1 _ {\mathrm{E}} = j ^ {*} \circ j = (p _ {1} \circ j) ^ {*} (p _ {1} \circ j) + b (u) ^ {*} b (u) = \mathrm{W} (u) + b (u) ^ {*} b (u),
\]

whence formula (3).

Let \(x \in \mathrm{dom}(u)\) and set \(y = \mathrm{W}(u)(x)\) . We have \(y \in \mathrm{dom}(u)\) and the formula \(y + u^{*}u(y) = x\) implies that \(u^{*}u(y) \in \mathrm{dom}(u)\) . We then have

\[
\begin{array}{r l} & {\mathrm{W} (u ^ {*}) ^ {- 1} (u (y)) = (1 _ {\mathrm{E}} + u u ^ {*}) (u (y))} \\ & {\qquad = u (y) + u u ^ {*} u (y) = u (y + u ^ {*} u (y)) = u (x),} \end{array}
\]

which means that the reduction of \(u \circ \mathrm{W}(u)\) to the domain of u is equal to \(\mathrm{W}(u^{*}) \circ u\) . Since the image of \(|p_{1}^{*}|\) is equal to that of \(p_{1}\), and hence to the domain of u (cor. of prop. 11 of I, p. 140), it follows

\[
u \circ \mathrm{W} (u) \circ | p _ {1} ^ {*} | = \mathrm{W} (u ^ {*}) \circ u \circ | p _ {1} ^ {*} | = \mathrm{W} (u ^ {*}) \circ b (u)
\]

(formula (2)). Moreover, \(|p_{1}^{*}| = \mathrm{W}(u)^{1/2}\) commutes with \(\mathrm{W}(u)\) , hence we obtain

\[
b (u) \circ \mathrm{W} (u) = u \circ | p _ {1} ^ {*} | \circ \mathrm{W} (u) = u \circ \mathrm{W} (u) \circ | p _ {1} ^ {*} | = \mathrm{W} (u ^ {*}) \circ b (u).
\]

This relation implies that \(b(u) \circ f(\mathrm{W}(u)) = f(\mathrm{W}(u^{*})) \circ b(u)\) for every function \(f \in \mathcal{C}(\mathbf{R}_{+})\) (prop. 11 of I, p. 113), hence in particular that \(b(u) \circ \mathrm{W}(u)^{1/2} = \mathrm{W}(u^{*})^{1/2} \circ b(u)\) . This proves formula (4) and concludes the proof.

COROLLARY. --- We have \(b(u)^{*} = b(u^{*})\) .

Denote by \(q_{1}\) and \(q_{2}\) the projections \(\Gamma_{u^{*}}\to\mathrm{E}\) and let \((k,|q_{1}^{*}|)\) be the polar decomposition of \(q_{1}^{*}\) , so that \(b(u^{*})=q_{2}\circ k\) . For all \(x\) and \(y\) in E, we have \(j(x)\in\Gamma_{u}\) and \(k(y)\in\Gamma_{u^{*}}\) , and since \(\Gamma_{u}\) is orthogonal to \(s(\Gamma_{u^{*}})\) by prop. 7 of IV, p. 236, it follows

\[
0 = \langle j (x) | s (k (y)) \rangle = \langle p _ {1} \circ j (x) | - q _ {2} \circ k (y) \rangle + \langle p _ {2} \circ j (x) | q _ {1} \circ k (y) \rangle .
\]

Since \(p_1 \circ j = |p_1^*|\) and \(q_1 \circ k = |q_1^*|\) (prop. 11, b) of I, p. 140) we obtain

\[
\langle b (u) x \mid | q _ {1} ^ {*} | y \rangle = \langle | p _ {1} ^ {*} | x \mid b (u ^ {*}) y \rangle ,
\]

hence \(|q_1^*|b(u) = b(u^*)^* |p_1^*|\) . Using formulas (1) and (4), we conclude from this that \(b(u)|p_1^*| = |q_1^*|b(u) = b(u^*)^* |p_1^*|\) , and since the image of \(|p_1^*|\) is the domain of \(u\) which is dense in E, we deduce \(b(u) = b(u^*)^*\) .

We denote \(\Omega(\mathrm{E})\) the set of \(v \in \mathcal{L}(\mathrm{E})\) such that \(1_{\mathrm{E}} - v^{*}v\) is positive and injective. For \(v \in \Omega(\mathrm{E})\) , the Hermitian endomorphism \((1_{\mathrm{E}} - v^{*}v)^{1/2}\) is injective, since its square is, and we denote by \(\mathrm{B}(v)\) the partial operator \(v \circ ((1_{\mathrm{E}} - v^{*}v)^{1/2})^{-1}\) .

Note that \(1_{\mathrm{E}} - v^{*}v\) is positive if and only if \(\| v\| \leqslant 1\) , since we have \(\langle x|(1_{\mathrm{E}} - v^{*}v)(x)\rangle = \| x\| ^2 -\| v(x)\| ^2\) for all \(x\in \mathrm{E}\) .

Lemma 1. --- The subset \(\Omega(\mathrm{E})\) is self-adjoint in \(\mathcal{L}(\mathrm{E})\) .

Let \(v \in \Omega(\mathrm{E})\) . We have \(\|v^{*}\| = \|v\| \leqslant 1\) . Therefore the endomorphism \(1_{\mathrm{E}} - vv^{*}\) is positive. It is injective: if \(x \in \operatorname{Ker}(1_{\mathrm{E}} - vv^{*})\) , we have \(vv^{*}(x) = x\) , whence \(v^{*}(v(v^{*}(x))) = v^{*}(x)\) , then \(v^{*}(x) = 0\) since \(1_{\mathrm{E}} - v^{*}v\) is injective, and finally \(x = v(v^{*}(x)) = 0\) . The lemma follows from this.

PROPOSITION 2. — The map \(u \mapsto b(u)\) defines a bijection from the set of closed densely defined partial operators on E onto the set \(\Omega(E)\) . The inverse bijection is given by \(v \mapsto B(v)\) .

The relation \(1_{\mathrm{E}} - b(u)^{*}b(u) = \mathrm{W}(u)\) (formula (3)) implies that \(b(u)\) belongs to \(\Omega (\mathrm{E})\) since \(\mathrm{W}(u)\) is positive and injective. Moreover, since \(b(u) = u\circ \mathrm{W}(u)^{1 / 2}\) (formula (2)), we obtain

\[
u = b (u) \circ (\mathrm{W} (u) ^ {1 / 2}) ^ {- 1} = b (u) \circ ((1 - b (u) ^ {*} b (u)) ^ {1 / 2}) ^ {- 1} = \mathrm{B} (b (u)).
\]

Conversely, let \(v \in \Omega(\mathrm{E})\) . Set \(w = (1_{\mathrm{E}} - v^{*}v)^{1/2}\) and set \(u = \mathrm{B}(v) = v \circ w^{-1}\) . The domain of u is the image of w, which is dense in E since w is Hermitian and injective (EVT, V, p. 41, prop. 4 (i)). For every \(x \in \mathrm{dom}(u)\) , it follows that

\[
\begin{array}{r l} & {\| u (x) \| ^ {2} = \langle (v ^ {*} v) (w ^ {- 1} (x)) | w ^ {- 1} (x) \rangle} \\ & {\qquad = - \langle (1 _ {\mathrm{E}} - v ^ {*} v) (w ^ {- 1} (x)) | w ^ {- 1} (x) \rangle + \| w ^ {- 1} (x) \| ^ {2}} \\ & {\qquad = - \langle w (x) | w ^ {- 1} (x) \rangle + \| w ^ {- 1} (x) \| ^ {2}} \end{array}
\]

since \((1_{\mathrm{E}} - v^{*}v) \circ w^{-1}\) is the reduction of w to the domain of \(w^{-1}\) . Since w is self-adjoint, it follows that

\[
\| w ^ {- 1} (x) \| ^ {2} = \| x \| ^ {2} + \| v \circ w ^ {- 1} (x) \| ^ {2} = \| x \| ^ {2} + \| u (x) \| ^ {2},
\]

and the partial operator \(\mathrm{B}(v) = u = v \circ w^{-1}\) is therefore closed by Lemma 2 of IV, p. 231.

Let us finally prove that \(b(\mathrm{B}(v)) = v\) . Since \(u = v \circ w^{-1}\) , we have \(v = u \circ w\) and it therefore suffices to check that \(w = \mathrm{W}(u)^{1/2}\) by formula (2), or even that \(1_{\mathrm{E}} - v^{*}v = \mathrm{W}(u)\) .

The endomorphism \(v^*\) belongs to \(\Omega(\mathrm{E})\) . By lemma 1, set \(w' = (1_{\mathrm{E}} - vv^{*})^{1/2}\) . It is an injective positive endomorphism of E, and we have \(v \circ w = w' \circ v\) (prop. 11 of I, p. 113). The graph of \(u\) is the set of elements of the form \((w(x), v(x))\) for \(x \in \mathrm{E}\) . By prop. 7 of IV, p. 236, the graph of \(u^*\) is \(s(\Gamma_u)^{\circ}\) . Since, for all \(x\) and \(y \in \mathrm{E}\) , we have

\[
\langle (w ^ {\prime} (x), v ^ {*} (x)) | (- v (y), w (y)) \rangle = - \langle x | w ^ {\prime} v (y) \rangle + \langle x | v w (y) \rangle = 0,
\]

the graph of \(u^*\) contains the elements of the form \((w'(x), v^*(x))\) for \(x \in \mathrm{E}\) , and we then have \(u^*(w'(x)) = v^*(x)\) . In particular, the domain of \(u^*\) contains the image of \(w'\) .

Let \(x \in E\) and set \(y = w^{2}(x) = (1_{\mathrm{E}} - v^{*}v)(x)\) , hence \(x = y + v^{*}v(x)\) . We have \(y \in \operatorname{dom}(u)\) and \(u(y) = v(w^{-1}(y)) = v(w(x)) = w'(v(x))\) . In particular, \(u(y) \in \operatorname{Im}(w) \subset \operatorname{dom}(u^{*})\) , and \(u^{*}(u(y)) = v^{*}(v(x))\) . It follows therefore \(x = y + v^{*}v(x) = y + u^{*}u(y)\) , whence \(y = \operatorname{W}(u)(x)\) , that is,

\[
(1 _ {\mathrm{E}} - v ^ {*} v) (x) = \mathrm{W} (u) (x).
\]

We have therefore shown that \(1_{\mathrm{E}} - v^{*}v = \mathrm{W}(u)\) , as desired.

COROLLARY. --- The endomorphism \(b(u)\) is Hermitian if and only if u is self-adjoint.

This follows from the injectivity of the map \(u \mapsto b(u)\) (prop. 2) and from the formula \(b(u^{*}) = b(u)^{*}\) (cor. of prop. 1).

